\documentclass[english, parskip=full]{scrartcl}
\usepackage[utf8]{inputenc}  % Kodierung der Datei
\usepackage[english]{babel}  % Sprache auf Deutsch 
\usepackage{color}
\usepackage{graphicx, gensymb}
\usepackage{amsfonts, cancel, amssymb}
\usepackage{amsmath, amsthm, array, esint}
\usepackage{booktabs, multirow}  % schönere Tabellen
\usepackage{caption}  % Anpassung der Captions
\usepackage{scrlayer-scrpage}    % Manipulation des Seitenstils
\pagestyle{scrheadings}  % Stil für die Seite setzen
\automark{section}  % Abschnittsnamen als Seitenbeschriftung 
\setheadsepline{.5pt}    % Kopzeile durch Linie abtrennen
\usepackage[hidelinks]{hyperref}  % Links und weitere PDF-Features
\usepackage{typearea, tikz, pgffor, verbatim}
\usetikzlibrary{calc, arrows.meta,arrows, graphs, quotes}
\usepackage[cache=false, outputdir=build]{minted}
\usemintedstyle{vs}
\usepackage{placeins} % provides \FloatBarrier

\definecolor{bg}{rgb}{0.9,0.9,0.9}
\newcommand\numberthis{\addtocounter{equation}{1}\tag{\theequation}}
\usepackage{svg}
\usepackage{siunitx}
\usepackage{physics}
\usepackage{tensor}
\renewcommand{\bra}{}
\newcommand{\kom}{}
\newcommand{\brak}{}
\renewcommand{\braket}{}
\newcommand{\intx}{}
\newcommand{\intr}{}
\newcommand{\kla}{}
\newcommand{\klam}{}
\newcommand{\klammer}{}
\newcommand{\oper}{}
\newcommand{\tecxt}{}
\newcommand{\Bra}{}
\newcommand{\Ket}{}
\renewcommand{\i}{\text{i}}
\newcommand{\e}{\text{e}}
\newcommand*{\Fourier}{\textsc{Fourier}\ }
\newcommand*{\F}{\mathcal{F}}
\newcommand*{\sinc}{\text{sinc\,}}
\newcommand{\conv}{\mathop{\scalebox{3.5}{\raisebox{-0.38ex}{\makebox[0.5\width][c]{$\ast$}}}}\limits}

\newcommand*{\titel}{3 Run and tumble particles in one dimension}
\newcommand*{\autor}{Joachim Schwardt}

\hypersetup{pdfauthor={\autor}, pdftitle={\titel}}

%\titlehead{titlehead}
\subject{Soft Condensed Matter}
\title{\titel}
\author{\autor}
\date{\today}

\begin{document}

%\begin{titlepage}
\maketitle  % Titel setzen
%\tableofcontents  % Inhaltsverzeichnis setzen
%\end{titlepage}

Consider particles in a finite one-dimensional region, that may be categorized as left-moving $L$ and right-moving $R$.
There is a fixed ``tumble-rate'' $\alpha$, with which each particle randomly changes its direction.
Let there be a symmetric potential $V(x) = \frac{\kappa}{2}x^2$.
This gives rise to a velocity $u := \frac{-V'(x)}{\gamma}$, where $\gamma>0$ is some constant.
We will also assume that the velocities in both directions may be different due to some unspecified reason.
Then the set of equations we need to solve is given by
\begin{align}
\partial_tR &= \alpha L - \alpha R - \partial_x\klam\left[ v_RR \right] - \partial_x\klam\left[ uR \right],\\
\partial_tL &= \alpha R - \alpha L - \partial_x\klam\left[ -v_LL \right] - \partial_x\klam\left[ uL \right].
\end{align}
Define the density $\rho := R+L$ and the polarization $\sigma := R-L$.
Also let $\bar{v} := \frac{v_R + v_L}{2}$ and $\Delta v := v_R - v_L$.
Then adding and subtracting the equations yields
\begin{align}
\partial_t\rho &= -\partial_x \klam\left[ u\rho \right] - \partial_x \klam\left[ \bar{v}\sigma + \frac{\Delta v}{2} \rho \right] \equiv -\partial_xJ(x),\\
\partial_t\sigma &= -2\alpha\sigma -\partial_x\klam\left[ u\sigma \right] - \partial_x\klam\left[ \bar{v}\rho + \frac{\Delta v}{2}\sigma \right].
\end{align}
The first equation has the form of a continuity equation, and thus the total density $\rho$ is a conserved quantity. 
This is also evident from the fact, that our model does not include a mechanism for removing or adding particles.
%The polarization is not conserved and decays on a timescale of $\tau \approx\frac{1}{2\alpha}$.
%We can show this by substituting $\sigma = \tilde{\sigma}\e^{-2\alpha t}$
%\begin{align}
%\partial_t\tilde{\sigma} &= -\partial_x\klam\left[ u\tilde{\sigma} \right] - \partial_x\klam\left[ \bar{v}\rho\e^{2\alpha t} + \frac{\Delta v}{2}\tilde{\sigma} \right].
%\end{align}
The polarization is not conserved due to the term $-2\alpha\sigma$ and thus decays on a timescale of $\tau \approx\frac{1}{2\alpha}$.
\\
In the steady state we have $\partial_t\rho=0=\partial_t\sigma$.
The continuity equation gives $\partial_xJ=0$ and thus 
\begin{align}
J&=u\rho + \frac{\Delta v}{2}\rho + \bar{v}\sigma = \tecxt\text{const} \equiv J_0.%,\\
%\Rightarrow\ \ \ \sigma &= \frac{J_0}{\bar{v}} - \frac{u + \frac{\Delta v}{2}}{\bar{v}} \rho.
\end{align}
Without the asymmetry induced by the $\Delta v$-term, we would expect $J_0=0$.
Therefore, we may assume that $J_0 \equiv \frac{\Delta v}{2}\rho$.
Then we find $\sigma = -\frac{u}{\bar{v}}\rho$.
Inserting this into the polarization equation yields
\begin{align}
0 &= 2\alpha \frac{u}{\bar{v}} \rho + \partial_x\klam\left[ \frac{\Delta v}{2} \frac{u}{\bar{v}}\rho \right] + \partial_x\klam\left[ \frac{u^2}{\bar{v}}\rho - \bar{v}\rho \right].
\end{align}
For our special case of $u = \frac{-\kappa x}{\gamma}$ this results in
\begin{align}
0 &= \frac{-2\alpha\kappa}{\gamma \bar{v}} x\rho - \partial_x\klam\left[ \frac{\kappa\Delta v }{2\gamma \bar{v}} x\rho \right] + \partial_x \klam\left[ \frac{\kappa^2 x^2}{\gamma^2\bar{v}} \rho - \bar{v}\rho \right] = \\
&= \frac{\kappa^2x^2}{\gamma^2\bar{v}} \rho'(x) - \bar{v}\rho' - \frac{\kappa\Delta v}{2\gamma\bar{v}} x\rho' + \frac{2\kappa^2}{\gamma^2\bar{v}} x\rho - \frac{\kappa \Delta v}{2\gamma\bar{v}} \rho - \frac{2\alpha\kappa}{\gamma\bar{v}} x\rho.
\end{align}
Cleaning up some of the coefficients, this is equivalent to
\begin{align}
0 &= \kappa x^2\rho' - \frac{\gamma\Delta v}{2} x\rho' - \frac{\gamma^2\bar{v}^2}{\kappa} \rho' + 2\klam\left[ \kappa - \alpha\gamma \right] x\rho - \frac{\gamma\Delta v}{2} \rho.
\end{align}
We may now separate
\begin{align}
\frac{\rho'(x)}{\rho(x)} &= \frac{2[\kappa -\alpha\gamma]x - \frac{\gamma\Delta v}{2}}{\frac{\gamma^2\bar{v}^2}{\kappa} + \frac{\gamma\Delta v}{2}x - \kappa x^2}.
\end{align}
Integrating gives
\begin{align}
\log(\rho) + c &= -\frac{\kappa - \alpha\gamma}{\kappa} \intx\int_{ }^{ } \frac{2x - \frac{\gamma\Delta v}{2[\kappa - \alpha\gamma]}}{x^2 - \frac{\gamma\Delta v}{2\kappa}x - \frac{\gamma^2\bar{v}^2}{\kappa^2}} \, \mathrm{d}x.
\end{align}
Let $x_{1,2} := -\frac{p}{2} \pm \sqrt{\frac{p^2}{4} - q}$.
Then $\frac{1}{(x-x_1)(x-x_2)} = \frac{A}{x-x_1} + \frac{B}{x-x_2}$ is a partial fraction decomposition.
We can find the coefficients by considering
\begin{align}
1 &= A(x-x_2) + B(x-x_1) = x(A+B) - (Ax_2 + Bx_1).
\end{align}
Thus $B=-A$ and $A=\frac{1}{x_1-x_2}$.
We now turn to solving the above integral in a general form
\begin{align}
I(x) &= \intx\int_{ }^{ } \frac{2x + a}{x^2 + px + q}\, \mathrm{d}x = \intx\int_{ }^{ } \frac{2x + p}{x^2+px+q} + \frac{a-p}{x^2+px+q} \, \mathrm{d}x = \\
&= \log\left\vert x^2+px+q\right\vert + (a-p) \intx\int_{ }^{ } \frac{1}{(x-x_1)(x-x_2)}\, \mathrm{d}x = \\
&= \log\left\vert x^2+px+q\right\vert + \frac{a-p}{x_1-x_2} \intx\int_{ }^{ } \frac{1}{x-x_1} - \frac{1}{x-x_2}\, \mathrm{d}x = \\
&= \log\left\vert (x-x_1)(x-x_2) \right\vert + \frac{a-p}{x_1-x_2} \log\left\vert \frac{x-x_1}{x-x_2} \right\vert = \\
&= \kla\left( 1 + \frac{a-p}{x_1-x_2} \right) \log |x-x_1| + \kla\left( 1 - \frac{a-p}{x_1-x_2} \right) \log |x-x_2|.
\end{align}
In our case we have $a = -\frac{\gamma\Delta v}{2[\kappa-\alpha \gamma]}$, $p=-\frac{\gamma\Delta v}{2\kappa}$ and $q=-\frac{\gamma^2\bar{v}^2}{\kappa^2}$.
The zeros are given by
\begin{align}
x_{1,2} &= \frac{\gamma\Delta v}{4\kappa} \pm \sqrt{\frac{\gamma^2\Delta v^2}{16\kappa^2} + \frac{\gamma^2\bar{v}^2}{\kappa^2}} = \frac{\gamma \Delta v}{4\kappa} \klam\left[ 1 \pm \sqrt{1 + \frac{16\bar{v}^2}{\Delta v^2}} \right].
\end{align}
Then the solution to the differential equation is
\begin{align}
\log (\rho) + c &= -\frac{\kappa - \alpha\gamma}{\kappa} \klam\left[ \log |x-x_1|^{r_1} + \log |x-x_2|^{r_2} \right],
\end{align}
where we defined the exponents 
\begin{align}
r_{1,2} &:= 1 \pm \frac{a-p}{x_1-x_2} = 1 \pm \frac{\frac{\gamma\Delta v}{2\kappa} - \frac{\gamma\Delta v}{2[\kappa - \alpha\gamma]}}{\frac{\gamma\Delta v}{2\kappa}\sqrt{1 + \kla\left( \frac{4\bar{v}}{\Delta v} \right)^2}} = 1 \pm \frac{1 - \frac{\kappa}{\kappa - \alpha\gamma}}{\sqrt{1 + \kla\left( \frac{4\bar{v}}{\Delta v} \right)^2}}.
\end{align}
The full solution is then given by
\begin{align}
\rho(x) &= c|x-x_1|^{-\frac{\kappa - \alpha\gamma}{\kappa}r_1} |x-x_2|^{-\frac{\kappa - \alpha\gamma}{\kappa}r_2},\\
\sigma(x) &= -\frac{u}{\bar{v}} \rho(x) = \frac{\kappa}{\gamma\bar{v}} x\rho(x).
\end{align}
Sanity check:
For the special case of $\Delta v=0$ we find $r_1=1=r_2$ and $x_{1,2} = \pm \frac{\gamma\bar{v}}{\kappa}$.
Then the solution takes the simple form
\begin{align}
\rho(x) &= \rho_0\left\vert \frac{x - x_1}{x - x_2} \right\vert^{\frac{\alpha\gamma}{\kappa} - 1} = \rho_0\left\vert \frac{x - \frac{\gamma\bar{v}}{\kappa}}{x + \frac{\gamma\bar{v}}{\kappa}} \right\vert^{\frac{\alpha\gamma}{\kappa} - 1},
\end{align}
which is (almost) identical to the result in the lecture.

\end{document}